\section{第二阶段深讲：进程生命周期不是 fork/exec/wait 三个词}

进程章节容易被讲成系统调用列表：\code{fork} 创建，\code{exec} 装载，\code{wait} 回收，\code{exit} 退出。这样的讲法太薄。真实内核维护的是一组互相引用的对象：任务、地址空间、文件表、信号状态、凭据、命名空间、调度实体、父子关系和资源统计。生命周期问题的核心是：什么时候对象对其他 CPU 可见，什么时候能失败，失败后怎样回滚，最后一个引用在哪里释放。

\subsection{从硬件执行流到 task}

硬件只知道当前 CPU 正在执行一组寄存器、PC、栈和页表根。操作系统把这组状态包装成 task。task 不等于进程的全部资源；它是可调度实体，指向更大的资源集合。

\begin{longtable}{p{0.2\linewidth}p{0.34\linewidth}p{0.34\linewidth}}
\toprule
资源 & 为什么不直接内嵌进 task & 生命周期问题 \\
\midrule
地址空间 & 线程可共享，exec 可替换 & 引用计数、COW、TLB shootdown \\
文件表 & fork 可共享，dup 共享 open file & close-on-exec、并发 close/read \\
信号处理 & 进程级和线程级状态交织 & handler、mask、pending queue \\
凭据 & 安全检查频繁读取，变更较少 & RCU 发布、setuid/exec 转换 \\
命名空间 & 容器内多个进程共享视图 & 引用计数、权限边界 \\
调度实体 & 每个线程都要被调度 & runqueue、vruntime、CPU affinity \\
\bottomrule
\end{longtable}

这种“task 指向对象”的结构让共享和替换成为可能。同一进程内多个线程共享 \code{mm} 和 \code{files}；\code{exec} 替换 \code{mm}，但保留 PID 和部分 fd；\code{fork} 复制 task，同时复制或共享不同对象。若把所有资源平铺在一个大结构里，这些语义会变得混乱。

\subsection{创建进程的真正顺序}

创建进程不能先把半成品放进 runqueue。只要任务进入 runnable 集合，另一个 CPU 就可能立刻运行它。因此 fork 的顺序必须从“不可见半成品”逐步推进到“对调度器可见”。

\begin{lstlisting}
copy_process outline:
  allocate task and kernel stack
  initialize minimal scheduler state
  copy credentials or install new references
  copy/share files, fs, signal state
  copy mm with COW page tables if needed
  allocate pid
  link child into parent relationships
  publish task to pid tables
  make task runnable only after invariants hold
\end{lstlisting}

每一步都有失败可能：内存不足、pid 不足、文件表复制失败、页表复制失败、权限检查失败。错误路径必须按逆序释放已经拿到的引用。这里的质量不在正常路径跑通，而在第 5 步失败时前 4 步是否完整回滚。

\begin{longtable}{p{0.22\linewidth}p{0.34\linewidth}p{0.32\linewidth}}
\toprule
阶段 & 成功后新增什么 & 失败时必须回滚什么 \\
\midrule
分配 task & task 内存和内核栈 & 释放栈和 task 对象 \\
复制 files & fd 表引用增加 & put 每个 open file 引用 \\
复制 mm & 页表页、PTE 引用、COW 标记 & 释放页表页、撤销引用 \\
分配 pid & pid 映射占用 & 释放 pid，不能留下僵尸映射 \\
链接父子 & 父进程 child list 可见 & 从列表摘除，避免 wait 看到半成品 \\
入队运行 & 调度器可选择任务 & 之后失败通常不能简单回滚，要走退出路径 \\
\bottomrule
\end{longtable}

这张表解释了为什么内核代码里有大量 \code{goto bad\_*}。这些标签表达的是阶段化回滚：从哪个阶段失败，就只回滚已经成功取得的资源。高质量内核代码宁愿显式写出阶段，也不能让半初始化对象泄漏到全局结构。

\subsection{fork 后父子为什么返回不同值}

fork 的一个经典语义是父进程返回子 PID，子进程返回 0。硬件没有“两个返回值”的指令，内核通过修改子任务保存的寄存器帧实现这个效果。父进程系统调用正常返回时，返回值寄存器写入 child pid；子任务第一次被调度时，从复制出的内核返回路径继续，返回值寄存器已经被设置为 0。

\begin{lstlisting}
parent path:
  ret = child.pid
  return_to_user(parent_regs with ret)

child prepared frame:
  child_regs = copy(parent_regs)
  child_regs.return_value = 0
  when scheduled, return_to_user(child_regs)
\end{lstlisting}

这个细节能帮助读者理解上下文切换：子进程不是从 \code{main} 开头重新执行，也不是内核调用了两次用户函数。它从和父进程相同的用户 PC 返回，只是寄存器状态略有不同。

\subsection{exec 是不可逆提交点}

\code{execve} 成功后，旧程序映像消失，新程序从入口点运行。难点在于失败语义：在提交点之前失败，应返回错误给旧程序；提交点之后失败，旧程序已经没有完整地址空间可恢复，只能终止任务。

\begin{lstlisting}
exec phases:
  copy argv/envp from old user memory
  open executable and check permissions
  parse binary format
  allocate new mm
  map text/data/interpreter/stack
  prepare credentials and auxv
  commit new mm and credentials
  jump to new entry
\end{lstlisting}

因此 exec 要尽量把可能失败的分配、权限检查和格式解析放在 commit 前。commit 动作要原子地切换 \code{mm}、设置新用户栈和入口、处理 close-on-exec fd、重置信号处理规则、应用 setuid/capability 变化。这个点之后，内核不能再假装旧程序还在。

\subsection{线程创建和资源共享边界}

线程和进程的差别不是“更轻量”四个字。线程共享地址空间、文件表和许多进程级状态，但拥有独立的可调度上下文、用户栈、内核栈、TLS、信号 mask 和调度状态。共享使通信便宜，也让错误传播更广。

\begin{longtable}{p{0.2\linewidth}p{0.34\linewidth}p{0.34\linewidth}}
\toprule
资源 & 进程间 fork 后 & 同进程线程间 \\
\midrule
地址空间 & 初始 COW，之后可分离 & 共享同一 \code{mm} \\
fd 表 & 可复制或共享，取决于语义 & 通常共享，close 影响其他线程 \\
用户栈 & 子进程有复制视图 & 每线程独立栈 \\
信号 mask & 每线程独立 mask & 进程 pending 与线程 pending 并存 \\
PID/TID & 父子不同 PID & 同 TGID，不同 TID \\
\bottomrule
\end{longtable}

线程库还要处理 TLS 和启动 trampoline。内核创建新线程时准备一个初始寄存器状态，使它第一次返回用户态时进入用户态线程入口，参数指向线程函数和参数对象。线程退出时要清理用户栈、TLS、robust futex、join 状态和内核 task 引用。

\subsection{exit 不是立刻消失}

任务调用 \code{exit} 后，不能简单释放自己。它可能还在别的表里有引用，父进程还要读取退出码，ptrace 或 wait 还可能观察它。于是退出分成多个阶段：停止用户执行、释放大资源、记录退出状态、通知父进程、进入 zombie，最后由 wait 或 reaper 释放 task。

\begin{lstlisting}
do_exit:
  mark task exiting
  close files and release mm
  release futexes and timers
  record exit_code and resource usage
  notify parent and reparent children
  become zombie
  schedule away forever

wait:
  find zombie child
  copy status to parent
  unlink child
  put final task reference
\end{lstlisting}

zombie 是一个必要状态：它保留父进程需要读取的结果，同时保证已退出任务不会再被调度。如果父进程不 wait，系统要通过 init 或 subreaper 接管，防止 zombie 永久积累。

\subsection{父子关系和 reparent}

进程树不是装饰。父进程负责 wait 子进程；信号传播、作业控制、会话、进程组都依赖关系结构。若父进程先退出，子进程不能失去回收者，内核会把它们 reparent 给 init 或最近的 subreaper。

\begin{lstlisting}
parent exits:
  for each child:
    choose new_parent = nearest subreaper or init
    move child to new_parent.children
    if child is zombie:
      notify new_parent
\end{lstlisting}

这里的并发点是父子列表、退出状态和 wait 扫描必须在锁保护下保持一致。否则 wait 可能漏掉刚退出的子进程，或者子进程被两个父对象同时看到。

\subsection{进程生命周期的测试矩阵}

进程生命周期测试应该覆盖正常路径、失败路径和并发路径。

\begin{longtable}{p{0.24\linewidth}p{0.42\linewidth}p{0.22\linewidth}}
\toprule
测试 & 关键断言 & 风险 \\
\midrule
fork 基本路径 & 父返回 child pid，子返回 0 & 寄存器帧设置错误 \\
fork 失败注入 & 每个阶段失败后引用数归零 & 资源泄漏或半成品可见 \\
fork COW & 父子写入互不污染 & PTE/TLB 权限错误 \\
exec 成功 & PID 保留，地址空间替换 & fd/signal/cred 规则错误 \\
exec 失败 & commit 前失败旧程序继续 & 旧 mm 被提前破坏 \\
exit/wait & zombie 被 wait 回收一次 & 重复回收或泄漏 \\
多线程 close/read & 共享 fd 表语义一致 & use-after-close 或引用泄漏 \\
reparent & 父退出后子由 reaper 回收 & zombie 泄漏 \\
\bottomrule
\end{longtable}

如果一本书只讲 fork/exec/wait 的定义，不讲这些断言，读者就无法写出可靠内核，也无法读懂真实源码里的复杂错误路径。
